% Options for packages loaded elsewhere
\PassOptionsToPackage{unicode}{hyperref}
\PassOptionsToPackage{hyphens}{url}
\documentclass[
  ignorenonframetext,
]{beamer}
\newif\ifbibliography
\usepackage{pgfpages}
\setbeamertemplate{caption}[numbered]
\setbeamertemplate{caption label separator}{: }
\setbeamercolor{caption name}{fg=normal text.fg}
\beamertemplatenavigationsymbolsempty
% remove section numbering
\setbeamertemplate{part page}{
  \centering
  \begin{beamercolorbox}[sep=16pt,center]{part title}
    \usebeamerfont{part title}\insertpart\par
  \end{beamercolorbox}
}
\setbeamertemplate{section page}{
  \centering
  \begin{beamercolorbox}[sep=12pt,center]{section title}
    \usebeamerfont{section title}\insertsection\par
  \end{beamercolorbox}
}
\setbeamertemplate{subsection page}{
  \centering
  \begin{beamercolorbox}[sep=8pt,center]{subsection title}
    \usebeamerfont{subsection title}\insertsubsection\par
  \end{beamercolorbox}
}
% Prevent slide breaks in the middle of a paragraph
\widowpenalties 1 10000
\raggedbottom
\AtBeginPart{
  \frame{\partpage}
}
\AtBeginSection{
  \ifbibliography
  \else
    \frame{\sectionpage}
  \fi
}
\AtBeginSubsection{
  \frame{\subsectionpage}
}
\usepackage{iftex}
\ifPDFTeX
  \usepackage[T1]{fontenc}
  \usepackage[utf8]{inputenc}
  \usepackage{textcomp} % provide euro and other symbols
\else % if luatex or xetex
  \usepackage{unicode-math} % this also loads fontspec
  \defaultfontfeatures{Scale=MatchLowercase}
  \defaultfontfeatures[\rmfamily]{Ligatures=TeX,Scale=1}
\fi
\usepackage{lmodern}
\ifPDFTeX\else
  % xetex/luatex font selection
\fi
% Use upquote if available, for straight quotes in verbatim environments
\IfFileExists{upquote.sty}{\usepackage{upquote}}{}
\IfFileExists{microtype.sty}{% use microtype if available
  \usepackage[]{microtype}
  \UseMicrotypeSet[protrusion]{basicmath} % disable protrusion for tt fonts
}{}
\makeatletter
\@ifundefined{KOMAClassName}{% if non-KOMA class
  \IfFileExists{parskip.sty}{%
    \usepackage{parskip}
  }{% else
    \setlength{\parindent}{0pt}
    \setlength{\parskip}{6pt plus 2pt minus 1pt}}
}{% if KOMA class
  \KOMAoptions{parskip=half}}
\makeatother
\usepackage{graphicx}
\makeatletter
\newsavebox\pandoc@box
\newcommand*\pandocbounded[1]{% scales image to fit in text height/width
  \sbox\pandoc@box{#1}%
  \Gscale@div\@tempa{\textheight}{\dimexpr\ht\pandoc@box+\dp\pandoc@box\relax}%
  \Gscale@div\@tempb{\linewidth}{\wd\pandoc@box}%
  \ifdim\@tempb\p@<\@tempa\p@\let\@tempa\@tempb\fi% select the smaller of both
  \ifdim\@tempa\p@<\p@\scalebox{\@tempa}{\usebox\pandoc@box}%
  \else\usebox{\pandoc@box}%
  \fi%
}
% Set default figure placement to htbp
\def\fps@figure{htbp}
\makeatother
\setlength{\emergencystretch}{3em} % prevent overfull lines
\providecommand{\tightlist}{%
  \setlength{\itemsep}{0pt}\setlength{\parskip}{0pt}}
\usepackage{bookmark}
\IfFileExists{xurl.sty}{\usepackage{xurl}}{} % add URL line breaks if available
\urlstyle{same}
\hypersetup{
  hidelinks,
  pdfcreator={LaTeX via pandoc}}

\author{\texorpdfstring{}{}}
\date{}

\begin{document}

\begin{frame}{Introduction: The Monetary Divide and the Architecture of
Value}
\protect\phantomsection\label{introduction-the-monetary-divide-and-the-architecture-of-value}
The late eighteenth century witnessed a fundamental restructuring of
British metaphysics and global monetary policy. Central to this upheaval
was the imposition of material quantification---an epistemological
framework derided by the poet, printmaker, and visionary William Blake:
``May God us keep / From Single vision \& Newtons sleep'' (Letter to
Thomas Butts, 22 November 1802) {[}@erdman1988complete{]}. Concurrently,
the British economic architecture was paralyzed by the meta-stable
tension of bimetallism. Gold and silver, theoretically bound by the
statutory fiat established by Isaac Newton on 21 September 1717
{[}@horsefield1960british{]}, operated not as cooperative halves of a
sovereign currency, but as autonomous vectors driven by international
arbitrage. Every monetary system must answer a foundational question:
\emph{what is the unit of account, and who controls its definition?} The
bimetallic answer---that two distinct physical substances could
simultaneously serve as a unified standard---was an inherent
contradiction, a legislative fiction sustained only so long as the
statutory ratio approximated fluctuating market realities.

\begin{block}{Gresham's Entropy and the Question of Value}
\protect\phantomsection\label{greshams-entropy-and-the-question-of-value}
This dynamic was the purest manifestation of Gresham's Law, explored
extensively in the subsequent section, the principle first articulated
by Thomas Gresham in 1558: when two metals circulate at a fixed ratio
misaligned with market values, the undervalued metal is melted and
exported while the overvalued metal remains {[}@thornton1802nature{]}.
In the British context, Newton's domestic ratio of 1:15.21 overvalued
gold against continental market ratios near 1:14.5, ensuring artificial
gold drove out undervalued silver. By 1715, nearly all the silver Newton
recoined during the Great Recoinage (1696--1699) had physically left the
country, shipped to France and the East Indies where its bullion weight
commanded a structural premium {[}@quinn1996competing;
@quinn2006bank{]}. The physical disappearance of metallic currency set
the necessary preconditions for the rise of abstract financial
instruments.
\end{block}

\begin{block}{Blake's Critique of Atomistic Commerce and Financial
Abstraction}
\protect\phantomsection\label{blakes-critique-of-atomistic-commerce-and-financial-abstraction}
Blake's lifelong critique of materialist atomism---the Epicurean
reduction of reality to disconnected, self-moving particles---maps
directly onto this escalating monetary abstraction
{[}@schouten2020epicurean{]}. While the Bank of England increasingly
relied on fractional reserves, paper issuance, and the arithmetic of
compounding imperial debt, Blake saw the fabric of human existence
divided into disconnected units of negotiable property. In \emph{London}
(Songs of Experience, 1794), the poet hears in ``every voice: in every
ban, / The mind-forg'd manacles'' (Plate 46, lines 7--8; Erdman p.~70)
{[}@erdman1988complete{]}---the invisible chains that commerce, law, and
state finance forge around the human imagination. As Mary Poovey
observes, the late eighteenth century saw the fundamental displacement
of intrinsic value by relational, speculative paper instruments---a
shift that alienated the laboring subject from their own work
{[}@poovey2008genres{]}.

Jon Mee notes that Blake's prophetic reaction to this crisis was
grounded in the radicalism of the 1790s, where political dissent fused
with critiques of state-sponsored economic policy
{[}@mee1992dangerous{]}. Blake's allegiance to this milieu is documented
in his annotations to Bishop Richard Watson's \emph{An Apology for the
Bible} (1798), defending Thomas Paine against establishment slander:
``To defend the Bible in this year 1798 would cost a man his life. The
Beast \& Whore damn the Bible'' (Erdman p.~616)
{[}@erdman1988complete{]}. E.P. Thompson extends this analysis, reading
Blake's deliberate imagery as a materialist indictment of the
acquisitive ethic---the systematic division of people and the
enforcement of moral bondage through market relations
{[}@thompson1993witness{]}. The ``dark Satanic Mills'' of the
\emph{Milton} preface (Erdman p.~95) operate simultaneously as the
literal industrial factories of the emerging capitalist order, and the
institutional ``mills'' of the established church grinding individuals
into doctrinal uniformity. As Joseph Albernaz identifies, measurement
technologies---whether the Newtonian standard of 1717 or the impending
1816 Gold Standard---enforce ``chimeric regimes.'' They attempt to
ground society in finite materiality while obfuscating humanity's shared
groundlessness {[}@albernaz2024common{]}.
\end{block}

\begin{block}{The Miser's Guinea and the Urizenic Inversion}
\protect\phantomsection\label{the-misers-guinea-and-the-urizenic-inversion}
This structural alienation extends from the macroeconomic sphere into
individual perception. Blake viewed the financial systems of his day not
merely as flawed economic policies, but as ontological regressions. In
his marginalia to Sir Joshua Reynolds' \emph{Discourses} (c.~1808),
Blake launches an explicit attack on this ideological inversion: ``To
the eyes of a miser a guinea is more beautiful than the sun, and a bag
worn with the use of money has more beautiful proportions than a vine
filled with grapes'' (Erdman p.~635) {[}@blake1808annotations{]}. The
golden guinea usurps the celestial sun; the dead container replaces the
living vine. This critique sharpens when Blake urgently demands in
\emph{Milton}: ``Leave counting Gold \& diamonds, for Gods eye delights
in thee'' (Plate 26, line 27; Erdman p.~121) {[}@erdman1988complete{]}.
For Blake, emergent capitalism and its central bankers represented a
conceptual flattening of existence, where the Enlightenment God
functioned as a senior accountant processing ledgers of abstract human
life rather than participating in its divine creation. The entire
structure of bimetallic generation and subsequent monometallic isolation
was a symptom of this Urizenic idolatry.
\end{block}

\begin{block}{Esoteric Finance and the Historiographic Framework}
\protect\phantomsection\label{esoteric-finance-and-the-historiographic-framework}
This historical backdrop has been documented by monetary historians like
Marc Flandreau and Alexander Del Mar.~As Flandreau argues, bimetallism
was an inherently unstable architecture, requiring a precarious
international equilibrium to prevent the collapse of domestic coinage
{[}@flandreau1996bimetallism{]}. Del Mar takes this further, viewing the
history of money as a constant struggle between state prerogative
(communal sovereignty) and private financial calculus (the
monopolization of coinage by mercantile elites)
{[}@delmar1895history{]}. David Erdman, in his foundational \emph{Blake:
Prophet Against Empire} (1954), connects Blake's cosmic visions directly
to political economy, interpreting Urizen as the ``great Work master''
who builds his cosmos through ``slave labor''---a direct textual
indictment of imperial war-financing and the commodification of human
energy {[}@erdman1954prophet{]}. Blake's prophetic works can be read as
an esoteric articulation of this struggle, a historically consistent
political myth that maps the spiritual descent into Ulro to the
materialist pressures of Bank Restriction-era credit economies
{[}@ferber2003politics; @duty2000sublime{]}.
\end{block}

\begin{block}{Manuscript Structure and Methodology}
\protect\phantomsection\label{manuscript-structure-and-methodology}
In this manuscript, we explicitly reconstruct the Newtonian Standard,
the 1797 Bank Restriction Act {[}@clapham1944bank;
@fetter1965development{]}, and the 1816 Coinage Act
{[}@redish1990evolution{]} not as rote economic touchstones, but as
topological fractures precipitating a systemic regression of the human
spirit. We attend to the \emph{minute particulars} of Blake's
engagement---his specific textual attacks on gold, his metallurgical
imagery, his precise plate compositions---while simultaneously tracing
the \emph{generic} structural logic by which any system of rigid
quantification must ultimately devour the qualitative life it claims to
measure. Blakean prophecy enacts a radical ideological inversion:
destroying the atomistic enclosures of Urizen, rejecting the spectral
flux of bimetallic generation, and laboring intensely at the Forges of
Los to rescue human agency and forge the ultimate ``Human Form Divine''
(detailed further under
\href{04a_illuminated_nonduality.md\#illuminated-nonduality-and-the-fourfold-vision}{Illuminated
Nonduality}).

\begin{figure}
\centering
\pandocbounded{\includegraphics[keepaspectratio,alt={Bimetallic and Prophetic Timeline --- Dual-Axis Chronology of Financial Collapse and Prophetic Ascent (1710--1900). Upper axis (Financial Events): progressive disintegration of the bimetallic synthesis from the 1717 Newtonian Standard (M: 1:15.21) through Hamilton's ratio of 15:1 (1792), the 1797 Bank Restriction, and the 1816 British Coinage Act (m\_S \textbackslash to 0), culminating in the ``Crime of 1873'' and Bryan's ``Cross of Gold'' (1896). Lower axis (Blake's Prophetic Works): Blake's inverse prophetic trajectory surges in metaphysical urgency---from The Book of Thel (1789) through America A Prophecy (1793), Songs of Experience (1794), The Four Zoas (1797--1807), and Jerusalem (1804--1820). Shaded era bands mark the British Bank Restriction Period (1797--1821) and the American Greenback Suspension (1862--1879)---twin clinamina of paper groundlessness on opposite shores of the Atlantic {[}@fetter1965development; @laughlin1898bimetallism{]}.}]{../output/figures/timeline.png}}
\caption{Bimetallic and Prophetic Timeline --- Dual-Axis Chronology of
Financial Collapse and Prophetic Ascent (1710--1900). Upper axis
(Financial Events): progressive disintegration of the bimetallic
synthesis from the 1717 Newtonian Standard (\(M: 1:15.21\)) through
Hamilton's ratio of 15:1 (1792), the 1797 Bank Restriction, and the 1816
British Coinage Act (\(m_S \to 0\)), culminating in the ``Crime of
1873'' and Bryan's ``Cross of Gold'' (1896). Lower axis (Blake's
Prophetic Works): Blake's inverse prophetic trajectory surges in
metaphysical urgency---from The Book of Thel (1789) through America A
Prophecy (1793), Songs of Experience (1794), The Four Zoas (1797--1807),
and Jerusalem (1804--1820). Shaded era bands mark the British Bank
Restriction Period (1797--1821) and the American Greenback Suspension
(1862--1879)---twin clinamina of paper groundlessness on opposite shores
of the Atlantic {[}@fetter1965development;
@laughlin1898bimetallism{]}.}\label{fig-timeline}
\end{figure}
\end{block}
\end{frame}

\end{document}
